%%
%% This is file `sample-sigconf.tex',
%% generated with the docstrip utility.
%%
%% The original source files were:
%%
%% samples.dtx  (with options: `all,proceedings,bibtex,sigconf')
%%
%% IMPORTANT NOTICE:
%%
%% For the copyright see the source file.
%%
%% Any modified versions of this file must be renamed
%% with new filenames distinct from sample-sigconf.tex.
%%
%% For distribution of the original source see the terms
%% for copying and modification in the file samples.dtx.
%%
%% This generated file may be distributed as long as the
%% original source files, as listed above, are part of the
%% same distribution. (The sources need not necessarily be
%% in the same archive or directory.)
%%
%%
%% Commands for TeXCount
%TC:macro \cite [option:text,text]
%TC:macro \citep [option:text,text]
%TC:macro \citet [option:text,text]
%TC:envir table 0 1
%TC:envir table* 0 1
%TC:envir tabular [ignore] word
%TC:envir displaymath 0 word
%TC:envir math 0 word
%TC:envir comment 0 0
%%
%% The first command in your LaTeX source must be the \documentclass
%% command.
%%
%% For submission and review of your manuscript please change the
%% command to \documentclass[manuscript, screen, review]{acmart}.
%%
%% When submitting camera ready or to TAPS, please change the command
%% to \documentclass[sigconf]{acmart} or whichever template is required
%% for your publication.
%%
%%
\documentclass[sigconf]{acmart}
%%
%% Get rid of copyright stuff
\setcopyright{none}
\settopmatter{printacmref=false}
\fancyhead{} % clear all header fields
\renewcommand\footnotetextcopyrightpermission[1]{} % Removes footnote with conference info
%%
%% Packages
\usepackage{listings} % File listings, with syntax highlighting
\usepackage{xcolor}   % Add coloring to syntax highlights of listings
\usepackage{ulem}     % For underlining
\usepackage{threeparttable}
%%
%%
%% Add Java listing config
\lstdefinestyle{java}{
    language=Java,
    basicstyle=\ttfamily\small,
    keywordstyle=\color{blue}\bfseries,
    stringstyle=\color{red},
    commentstyle=\color{gray}\itshape,
    numbers=left,
    numberstyle=\tiny,
    stepnumber=1,
    numbersep=10pt,
    showspaces=false,
    showstringspaces=false,
    showtabs=false,
    frame=single,
    tabsize=2,
    captionpos=b,
    breaklines=true,
    breakatwhitespace=false,
    escapeinside={\%*}{*)},
    morecomment=[l][\color{gray}]{//},
    morecomment=[s]{/*}{*/},
    morestring=[b]",
    morekeywords={public, class, private, protected, void, int, if, else, return, true, false, static, new}
}
%%
%%
%% Add C++ listing config
\lstdefinestyle{cpp}{
    language=C++,
    basicstyle=\ttfamily\small,
    keywordstyle=\color{blue},
    commentstyle=\color{gray},
    captionpos=b,
    numbers=left,
    numberstyle=\tiny,
    breaklines=true,
    frame=single,
}
%%
%%
%% Add Python listing config
\lstdefinestyle{python}{
    language=Python,
    basicstyle=\ttfamily\small,
    keywordstyle=\color{purple},
    commentstyle=\color{gray},
    captionpos=b,
    numbers=left,
    numberstyle=\tiny,
    breaklines=true,
    frame=single,
}
%%
%%
%% \BibTeX command to typeset BibTeX logo in the docs
\AtBeginDocument{%
  \providecommand\BibTeX{{%
    Bib\TeX}}}
%% end of the preamble, start of the body of the document source.
\begin{document}

%%
%% The "title" command has an optional parameter,
%% allowing the author to define a "short title" to be used in page headers.
\title{LISP Interpreter Report}

%%
%% The "author" command and its associated commands are used to define
%% the authors and their affiliations.
%% Of note is the shared affiliation of the first two authors, and the
%% "authornote" and "authornotemark" commands
%% used to denote shared contribution to the research.
\author{Amairany Neri}
\affiliation{
  \institution{New Mexico Institute of Mining and Technology Deparment of Computer Science}
  \city{Socorro}
  \state{NM}
  \country{USA}
}
\email{amairany.neri@student.nmt.edu}

%%
%% By default, the full list of authors will be used in the page
%% headers. Often, this list is too long, and will overlap
%% other information printed in the page headers. This command allows
%% the author to define a more concise list
%% of authors' names for this purpose.
\renewcommand{\shortauthors}{Neri}
\acmConference[]{}{}{}

%%
%% The abstract is a short summary of the work to be presented in the
%% article.
\begin{abstract}
    In this paper, we discuss the implementation of an interactive LISP interpreter using Python.
    The main goal of this project was to understand the concepts behind creating an interpreter
    for the simple functional programming language LISP.
\end{abstract}

%%
%% The code below is generated by the tool at http://dl.acm.org/ccs.cfm.
%% Please copy and paste the code instead of the example below.
%%
\begin{CCSXML}
<ccs2012>
<concept>
<concept_id>10002944</concept_id>
<concept_desc>General and reference</concept_desc>
<concept_significance>500</concept_significance>
</concept>
</ccs2012>
\end{CCSXML}

\ccsdesc[500]{General and reference}

%%
%% Keywords. The author(s) should pick words that accurately describe
%% the work being presented. Separate the keywords with commas.
\keywords{LISP, Interpreter, Functional Programming, Python}
%% This command processes the author and affiliation and title
%% information and builds the first part of the formatted document.
\maketitle

% Introduction
\section{Introduction}
The LISP Interpreter was created using Python. Very similar to the
actual LISP interpreter, you input LISP expressions and the interpreter
evaluates it on the line below. An output file of the session is generated
under the name \texttt{results.file}.

% Project Final Status Reporting
\section{LISP Constructs Implemented}
Below is a table of the status of each LISP construct implemented in the interpreter.
\begin{table}[h!]
\begin{threeparttable}
    \caption{Project Progress/Final Status Reporting}
    \begin{tabular}{| c | c |}
    \hline
    Lisp Construct & Status \\
    \hline
    Variable reference & DONE \\
    Constant literal & DONE \\
    Quotation & DONE \\
    Conditional & DONE \\
    Variable Definitions & DONE \\
    Functional call & DONE \\
    Assignment & DONE \\
    Function Definition & DONE \\
    The arthimetic +, -, *, / operators on INTEGER type & DONE \\
    "car" and "cdr & DONE \\
    The built-in function: "cons" & DONE \\
    sqrt, exp & DONE \\
    >, <, =, <=, >=, != & DONE \\
    \hline
    \end{tabular}
    \begin{tablenotes}
      \item[1] \texttt{set!} lacks error checking for formal parameters.
      \item[2] Division uses floating-point arithmetic.
    \end{tablenotes}
    \label{tab:lispconstructs}
\end{threeparttable}
\end{table}

The following functions for extra credit were also implemented:
\begin{table}[h!]
    \begin{tabular}{| c | c|}
        \hline
        Extra credits & Status \\
        \hline
        "mapcar" & DONE \\
        "lambda" & DONE \\
        \hline
    \end{tabular}
    \label{tab:extrafunctions}
    \caption{Extra Credits Functions}
\end{table}

\section{Implementation Descriptions}
The interpreter implements core Lisp constructs through a combination of recursive parsing and environment-based evaluation \cite{lispy}\cite{pythondoc}. Variable references are resolved via the \texttt{Environment.find()} method, which recursively searches parent scopes until a match is found. Quotation is handled during parsing by wrapping quoted expressions in a \texttt{(quote ...)} list that bypasses evaluation. For conditionals, the \texttt{if} special form evaluates the test expression first, then selectively evaluates either the consequent or alternative branch using Python's ternary logic.
Function calls leverage Python's first-class function capabilities. Built-in operators like \texttt{+} and \texttt{car} are stored as lambda functions in the global environment dictionary, while user-defined functions created via \texttt{defun} or \texttt{lambda} are instantiated as \texttt{Procedure} objects. These objects capture their defining environment during creation to enable static scoping, ensuring proper variable resolution during subsequent calls.
\section{Applied Class Concepts}
The implementation directly applies several core concepts from programming language theory. Static scoping is achieved through environment chaining, where each \texttt{Procedure} maintains a reference to its birth environment. Recursive descent parsing techniques learned during syntax analysis modules were adapted to convert token streams into abstract syntax trees. The interpreter's memory model mirrors stack-based execution through nested \texttt{Environment} objects, though without explicit stack frames.

First-class functions are implemented via the \texttt{Procedure} class wrapping, demonstrating higher-order function principles. Error handling for arithmetic operations incorporates concepts from exception handling lectures, with overflow checks using 32-bit integer bounds. The \texttt{mapcar} implementation utilizes list comprehension and dynamic function application, reflecting functional programming paradigms discussed in class.

\section{Implementation Details}
The interpreter employs several notable design choices. Boolean values \texttt{T} and \texttt{NIL} are mapped to Python's native \texttt{True} and \texttt{False}, with special formatting during output conversion. The output file mechanism uses Python's context manager pattern, ensuring proper file closure even during error conditions.

Environment lookups use chain-of-responsibility pattern through the \texttt{find()} method rather than prototype inheritance. Arithmetic operators are wrapped in safety functions (\texttt{safe\_add}, etc.) to enforce 32-bit integer constraints while maintaining Python's native number types. The \texttt{cons} implementation automatically flattens improper lists to match Lisp's expected behavior.

\section{Conclusion}
This project provided significant insights into both Lisp semantics and interpreter design. Implementing special forms like \texttt{quote} and \texttt{if} deepened understanding of evaluation strategies, while environment management illuminated scoping concepts. Challenges in balancing Python's dynamic typing with Lisp's type expectations highlighted the importance of careful edge case handling.

The experience reinforced the elegance of Lisp's minimal syntax paired with powerful abstraction capabilities. Future improvements could add proper tail recursion and expand the type system, but the current implementation successfully demonstrates core interpreter concepts while adhering to project requirements.

%%
%% The next two lines define the bibliography style to be used, and
%% the bibliography file.
\bibliographystyle{ACM-Reference-Format}
\bibliography{project-references}


\end{document}
\endinput
%%
